% Generated by tools/edn_latex.py from reports/edn.md.
% Do not edit: change reports/edn.md and run `make edn`.
\ednentry{
  id = {EDN-53},
  title = {\texorpdfstring{\texttt{train.\allowbreak{}num\_\allowbreak{}threads}}{train.num\_threads} is pinned to 1, trading fit speed for a machine-independent artefact},
  date = {2026-09-30},
  milestone = {M3: Quality Assurance},
  topic = {Reproducibility, Modelling Approach},
  participants = {@lukas2510, decided by the repository owner},
  decision = {\texttt{params.\allowbreak{}yaml} carries \texttt{train.\allowbreak{}num\_\allowbreak{}threads: 1} and \texttt{train} passes it to LightGBM as \texttt{n\_\allowbreak{}jobs}. It is never left at LightGBM's default of \texttt{-\allowbreak{}1}, which means every core. The reproducibility it buys is a \texttt{booster.\allowbreak{}txt} that is byte-identical on any machine. What it costs is narrower than first recorded: measured against a pin of 2 or 4 threads it costs about 1.5x the fit, and measured against the default it actually \emph{saves} time and, more to the point, removes an enormous variance.},
  rationale = {\emph{Alternatives considered:}
\begin{itemize}[leftmargin=*, topsep=2pt]
\item \textbf{Option A: leave \texttt{n\_\allowbreak{}jobs} at its default of \texttt{-\allowbreak{}1}.}
\newline Pros: nothing to configure. Not, as it turns out, the fastest fit: see the rationale.
\newline Cons: LightGBM writes the thread count into \texttt{booster.\allowbreak{}txt} as a \texttt{[num\_\allowbreak{}threads: N]} line, so the artefact -- and therefore its DVC hash and \texttt{dvc.\allowbreak{}lock} -- depends on the core count of whoever ran \texttt{dvc repro}. Two contributors with different laptops would produce different artefacts from identical data and parameters, and \texttt{dvc status} could never say the model was unchanged. On top of that, on an 8-core machine the default is both slower and wildly variable at real scale, and it makes the test suite's cost unpredictable, which matters because the suite fits all four variants on 744 rows.
\item \textbf{Option B (chosen): pin to 1.}
\newline Pros: makes \texttt{booster.\allowbreak{}txt} byte-identical across thread counts, which is a stronger claim than NFR-06's ``metrics within 0.1 percentage points'' and is the one the reproducibility test asserts. It also removes oversubscription if the four independent \texttt{train@\allowbreak{}*} stages are ever run at once, and it is the value with by far the smallest spread between repeated fits at either data scale -- which is what makes a reported fit time worth reporting.
\newline Cons: it is measurably slower than a \textbf{pinned} 2 or 4. On the 60,378-row training split \texttt{lgbm-\allowbreak{}extended} fits in a median 18.3 s at 1 thread against 12.2 s at 4, so option C would buy back about a third of the fit per variant. It is not slower than the default it actually replaces.
\item \textbf{Option C: pin to a fixed larger number, 2 or 4.}
\newline Pros: the same machine independence as option B, because the value is pinned rather than discovered, with most of the speed back. This, and not option A, is the alternative that the pin genuinely costs something against.
\newline Cons: it is only machine-independent on machines that \emph{have} that many cores. LightGBM silently uses what is available, so a 2-core CI runner asked for 4 would write a different \texttt{booster.\allowbreak{}txt} than an 8-core laptop, and the property would hold right up until the first machine that breaks it -- the worst kind of guarantee. 1 is the only value every machine can honour.
\end{itemize}
\par
\emph{Why:} The measurement that decides the direction is that the \textbf{trees are already thread-independent and only the recorded setting is not}. Fitting \texttt{lgbm-\allowbreak{}basic} on the fixture in separate processes at 1, 2, 4 and 8 threads gives the same \texttt{best\_\allowbreak{}iteration} (98) and a byte-identical model file once the \texttt{[num\_\allowbreak{}threads: N]} line is normalised away, at every thread count; the md5 of the whole file differs at every one of them. So pinning the value does not make the model reproducible -- it already is -- it makes the \emph{file} reproducible, which is what DVC hashes.
\par
That also corrects a claim worth not repeating. The design this work followed said \texttt{deterministic=\allowbreak{}True} and \texttt{force\_\allowbreak{}row\_\allowbreak{}wise=\allowbreak{}True} are what keep the trees stable across thread counts. They are not, at this data size: with both flags \textbf{off}, the normalised file is still byte-identical across 1, 2, 4 and 8 threads. What the flags do change is \emph{which} trees are built -- the normalised md5 differs between the flags on and off -- so they are not free of consequence, they are simply not what makes the result thread-independent here. They are kept as insurance documented by LightGBM, and the code and the model card now say that rather than claiming a measurement that does not exist.
\par
The cost side was measured rather than assumed, and the first version of this entry stated it against a configuration the default never produces. The full real-scale curve, medians of three fits per thread count on the 60,378-row training split, on an 8-core machine with nothing else running:
\par
{\small\renewcommand{\arraystretch}{1.15}%
\begin{tabularx}{\linewidth}{@{}>{\raggedright\arraybackslash}X>{\raggedright\arraybackslash}X>{\raggedright\arraybackslash}X>{\raggedright\arraybackslash}X>{\raggedright\arraybackslash}X@{}}
\toprule
Variant & 1 thread & 2 & 4 & 8 (the default here) \\
\midrule
\texttt{lgbm-\allowbreak{}basic} & 7.1 s (7.0 to 7.2) & 5.6 s (5.3 to 6.8) & 6.8 s (6.4 to 7.4) & 17.7 s (16.5 to 36.5) \\
\texttt{lgbm-\allowbreak{}extended} & 18.3 s (17.9 to 18.9) & 12.8 s (12.7 to 13.1) & 12.2 s (11.3 to 14.3) & 29.2 s (12.6 to 38.7) \\
\bottomrule
\end{tabularx}}
\par
Two things follow, and the second is the correction. More threads do help, so the design's ``there is no performance reason to use more'' is wrong: 2 to 4 threads are consistently the fastest, and the pin costs about 1.5x the fit for \texttt{lgbm-\allowbreak{}extended} and 1.27x for \texttt{lgbm-\allowbreak{}basic} against them. But the alternative this decision actually rejects is not a pin of 2 or 4, it is LightGBM's \textbf{default of every core}, and against that the pin is free or a gain: 18.3 s against 29.2 s for \texttt{lgbm-\allowbreak{}extended}, 7.1 s against 17.7 s for \texttt{lgbm-\allowbreak{}basic}. So the trade is with option C, not with option A.
\par
The spread matters more than the median. At 1 thread the three fits of a variant land within 0.2 s and 1.0 s of each other; at 8 they span 20 s and 26 s, and the fastest \texttt{lgbm-\allowbreak{}extended} fit of the whole sweep (12.6 s) and the slowest (38.7 s) are both at 8 threads. A number that varies threefold between identical runs is not a fit time anyone can report, which is a reproducibility argument for the pin quite apart from the artefact's hash.
\par
The earlier entry attributed the reversal at 8 threads to the 744-row fixture, where the direction is even sharper: 0.33 s at 1 thread, 0.27 s at 2, 0.48 s at 4 and a median of \textbf{28.58 s at 8}, with one fit of 143.71 s, on a machine already at load 13 to 19. That reading was too narrow. The reversal is not about the row count but about the work per thread: it happens at 60,378 rows too, and on \texttt{lgbm-\allowbreak{}basic} -- 16 features -- more decisively than on \texttt{lgbm-\allowbreak{}extended}'s 162, which is where the extra threads have something to divide. Since the test suite fits every variant on the fixture, on whatever CI runner picks up the job, the pin is what keeps the suite's cost predictable as well.
\par
Absolute seconds here are machine- and load-dependent and are not comparable with the ladder table's one-off fit times in the model card; the ratios within one sweep are the measurement.},
  ai = {Information seeking, Alternative assessment, Recommendation, Solution generation},
  response = {Accepted with modifications},
  assessment = {AI identified that LightGBM records the thread count in the artefact, which is the fact the whole decision rests on, and proposed pinning it. Three of its supporting claims were wrong and were corrected by measuring. That \texttt{deterministic=\allowbreak{}True} and \texttt{force\_\allowbreak{}row\_\allowbreak{}wise=\allowbreak{}True} are what make the trees thread-independent: they are not at this size, the trees are identical with the flags off, though the flags do change which trees are built. That pinning to 1 costs nothing at real scale: it costs about 1.5x the fit against a pin of 2 to 4. And then the correction to that correction -- the entry's own second version stated the cost against 2 to 4 as if that were what the default does, and put eight threads at real scale under ``not measured''. Measuring it reversed the comparison that matters: against the default the pin is free or a gain. The entry now names which alternative the cost is against, and the code comments were rewritten to match.},
  aievidence = {Claude Code session on 2026-09-30 implementing issue \#37: \texttt{booster.\allowbreak{}txt} was hashed with the \texttt{[num\_\allowbreak{}threads: N]} line normalised away across four thread counts and both settings of the determinism flags. The thread sweep was then redone during the review of pull request \#62, as medians of three fits at 1, 2, 4 and 8 threads on both LightGBM variants at real scale on an idle machine, after a first attempt was discarded for having shared the machine with a test run -- which is also how the size of the variance at 8 threads came to light.},
  evidence = {\href{https://github.com/mlops-2627q1-mds-upc/MLOPS_Recommenditos/blob/main/recommenditos/modeling/model.py}{\texttt{recommenditos/\allowbreak{}modeling/\allowbreak{}model.\allowbreak{}py}}, \texttt{LightGBMModel.\allowbreak{}fit}; the \texttt{num\_\allowbreak{}threads} comment in \href{https://github.com/mlops-2627q1-mds-upc/MLOPS_Recommenditos/blob/main/params.yaml}{\texttt{params.\allowbreak{}yaml}}; \texttt{tests/\allowbreak{}test\_\allowbreak{}model.\allowbreak{}py::test\_\allowbreak{}num\_\allowbreak{}threads\_\allowbreak{}is\_\allowbreak{}pinned\_\allowbreak{}from\_\allowbreak{}params}, which trains at a value the project does not use so a hard-coded 1 would fail it, and \texttt{::test\_\allowbreak{}the\_\allowbreak{}payload\_\allowbreak{}of\_\allowbreak{}a\_\allowbreak{}refit\_\allowbreak{}is\_\allowbreak{}byte\_\allowbreak{}identical}; \href{https://github.com/mlops-2627q1-mds-upc/MLOPS_Recommenditos/blob/main/docs/docs/model-card.md}{model card}, Training Procedure, Determinism; \href{https://github.com/mlops-2627q1-mds-upc/MLOPS_Recommenditos/blob/main/docs/docs/requirements.md}{requirements} NFR-06 and NFR-10; \href{https://github.com/mlops-2627q1-mds-upc/MLOPS_Recommenditos/issues/37}{issue \#37}.},
}
